\section{Conclusion}
\label{sec:conclusion}

\trustepisode provides a deterministic, auditable EDR/NDR episode-reconstruction contract:
single-owner immutable revisions, append-only terminal outcomes with a complete failure-code
matrix, and stage-localized A--G evidence-loss diagnosis. On the sealed M1--M6 cohort, sliding
sessionization improves completeness relative to fixed bins, but fair sliding time/entity/relation
builders show \emph{no observed difference} from full TrustEpisode (60/60 ties). That negative
result is forensic evidence about benchmark coverage: advanced guards require targeted
conformance cards, which we provide without treating them as live effectiveness. Canonical
replay is byte-identical on all 60 sealed malicious runs; raw arrival-order robustness requires
re-canonicalization. M2's missing \texttt{file\_transfer} token is
\texttt{D\_token\_mapping\_failure}, not a demonstrated formation drop. Contract ablation shows
which properties prevent mutable history, silent missing-input failure, and noncanonical replay.
Supplemental partition artifacts (not primary estimators) include L1 live bases with offline
L1/L2 ablations, L3 controlled hub inject, L4 controlled late inject, wall-clock within-horizon
holds (2/2), and sealed beyond-horizon attempts without organic LateEvidenceRecord. Sealed M7 is
a supplemental 3900\,s boundary-stress cohort only. Fitted scoring, calibration, and
telemetry-outage robustness remain unevaluated and outside these claims.
